\SixthArticle{27}
Das Integral $\int x^{\lambda-1}(1-x^\mu)^\nu\,\mathrm{d}x$, so gewählt, dass es für $x=0$ verschwindet, wird durch die folgende Reihe ausgedrückt, sofern $\lambda,\mu$ positive Größen sind:
\[
\frac{x^\lambda}{\lambda}
-\frac{\nu x^{\mu+\lambda}}{\mu+\lambda}
+\frac{\nu(\nu-1)x^{2\mu+\lambda}}{1.2.(2\mu+\lambda)}
-\mathrm{usw.}
=\frac{x^\lambda}{\lambda}F(-\nu,\frac\lambda\mu,
\frac\lambda\mu+1,x^\mu)
\]
Sein Wert für $x=1$ ist daher
\[
=\frac{\Pi\frac\lambda\mu.\Pi\nu}
{\lambda\Pi(\frac\lambda\mu+\nu)}.
\]
Aus diesem Satz ergeben sich von selbst alle Beziehungen, die der berühmte \textsc{Euler} einst mit großer Mühe entwickelt hat. Setzt man beispielsweise
\[
\int\frac{\mathrm{d}x}{\BareRoot{(1-x^4)}}=A,\qquad
\int\frac{xx\,\mathrm{d}x}{\BareRoot{(1-x^4)}}=B
\]
so ergibt sich $A=\frac{\Pi\frac14.\Pi(-\frac12)}{\Pi(-\frac14)}$, $B=\frac{\Pi\frac34.\Pi(-\frac12)}{3\Pi\frac14}
=\frac{\Pi(-\frac14).\Pi(-\frac12)}{4\Pi\frac14}$ und somit $AB=\frac14\pi$. Zugleich folgt daraus, da $\Pi\frac14.\Pi(-\frac14)=\frac{\frac14\pi}{\sin\frac14\pi}
=\frac\pi{\BareRoot{8}}$ ist,
\[
\Pi\frac14=\BareRoot[4]{(\frac18\pi AA)}
=\BareRoot[4]{\frac{\pi^3}{128BB}},\qquad
\Pi(-\frac14)=\BareRoot[4]{\frac{\pi^3}{8AA}}
=\BareRoot[4]{(2\pi BB)}
\]
Der Zahlenwert von $A$ ist nach der Berechnung von \textsc{Stirling} $=1{,}3110287771\;\allowbreak 4605987$; der Wert von $B$ ist nach demselben Autor $=0{,}5990701173\;\allowbreak 6779611$, nach unserer Berechnung, die auf einem besonderen Kunstgriff beruht, $=0{,}5990701173\;\allowbreak 6779610372$.

Allgemein lässt sich leicht zeigen, dass der Wert der Funktion $\Pi z$, wenn $z$ eine rationale Größe $=\frac m\mu$ ist und $m,\mu$ ganze Zahlen bezeichnen, aus $\mu-1$ bestimmten Werten solcher Integrale für $x=1$ hergeleitet werden kann, und zwar auf sehr viele verschiedene Arten. Nimmt man nämlich für $\lambda$ eine ganze Zahl und für $\nu$ einen Bruch, dessen Nenner $=\mu$ ist, so lässt sich der Wert jenes Integrals stets auf drei $\Pi z$ zurückführen, wobei $z$ ein Bruch mit dem Nenner $=\mu$ ist; jedes derartige $\Pi z$ lässt sich aber nach Formel 45 entweder auf $\Pi(-\frac1\mu)$ oder auf $\Pi(-\frac2\mu)$ oder auf $\Pi(-\frac3\mu)$ usw. oder auf $\Pi(-\frac{\mu-1}{\mu})$ zurückführen, sofern $z$ tatsächlich ein Bruch ist; denn ist $z$ eine ganze Zahl, so ist $\Pi z$ an sich bekannt. Aus jenen Werten der Integrale lässt sich aber, allgemein gesprochen, jedes $\Pi(-\frac m\mu)$,
\SourceLeaf{160}
wenn $m<\mu$ ist, durch Elimination gewinnen\HistoricalNote{17}.
Sogar halb so viele derartige Integrale genügen, wenn wir zugleich Formel 54 zu Hilfe nehmen. Setzt man beispielsweise
\[
\begin{gathered}
\int\frac{\mathrm{d}x}{\BareRoot[5]{(1-x^5)}}=C,\quad
\int\frac{\mathrm{d}x}{\BareRoot[5]{(1-x^5)^2}}=D,\quad
\\\int\frac{\mathrm{d}x}{\BareRoot[5]{(1-x^5)^3}}=E,\quad
\int\frac{\mathrm{d}x}{\BareRoot[5]{(1-x^5)^4}}=F,
\end{gathered}
\]
so ergibt sich
\[
\begin{gathered}
C=\Pi\frac15.\Pi(-\frac15),\quad
D=\frac{\Pi\frac15.\Pi(-\frac25)}{\Pi(-\frac15)},\quad
\\E=\frac{\Pi\frac15.\Pi(-\frac35)}{\Pi(-\frac25)},\quad
F=\frac{\Pi\frac15.\Pi(-\frac45)}{\Pi(-\frac35)}
\end{gathered}
\]
Da $\Pi\frac15=\frac15\Pi(-\frac45)$ ist, haben wir daher
\[
\begin{gathered}
\Pi(-\frac15)=\BareRoot[5]{\frac{5C^4}{DEF}},\quad
\Pi(-\frac25)=\BareRoot[5]{\frac{25C^3D^3}{EEFF}},\quad
\\\Pi(-\frac35)=\BareRoot[5]{\frac{125CCDDEE}{F^3}},\\
\Pi(-\frac45)=\BareRoot[5]{(625CDEF)}
\end{gathered}
\]
Die Formeln 54 und 55 ergeben ferner
\[
C=\frac\pi{\sin\frac15\pi},\qquad
\frac DF=\frac{\sin\frac15\pi}{\sin\frac25\pi}
\]
sodass die beiden Integrale $D,E$ oder $E$ und $F$ genügen, um alle Werte $\Pi(-\frac15)$, $\Pi(-\frac25)$ usw. zu berechnen.

\SixthArticle{28}
Setzt man $y=\nu x$ und $\mu=1$, so ist $\frac{\Pi\lambda.\Pi\nu}{\lambda\Pi(\lambda+\nu)}$ der Wert des Integrals $\int\frac{y^{\lambda-1}(1-\frac y\nu)^\nu\,\mathrm{d}y}{\nu^\lambda}$ von $y=0$ bis $y=\nu$, oder der Wert des Integrals $\int y^{\lambda-1}(1-\frac y\nu)^\nu\,\mathrm{d}y$ zwischen denselben Grenzen ist $=\frac{\nu^\lambda\Pi\lambda.\Pi\nu}{\lambda\Pi(\lambda+\nu)}
=\frac{\Pi(\nu,\lambda)}\lambda$ (Form. 47), sofern $\nu$ eine ganze Zahl bezeichnet. Wächst nun $\nu$ ins Unendliche, so ist der Grenzwert von $\Pi(\nu,\lambda)$ $=\Pi\lambda$, der Grenzwert von $(1-\frac y\nu)^\nu$ dagegen $e^{-y}$, wobei $e$ die Basis der natürlichen Logarithmen bezeichnet. Wenn also $\lambda$ positiv ist, drückt $\frac{\Pi\lambda}\lambda$ oder $\Pi(\lambda-1)$ das Integral $\int y^{\lambda-1}e^{-y}\,\mathrm{d}y$ von $y=0$ bis $y=\infty$ aus; oder, wenn $\lambda$ für $\lambda-1$ geschrieben wird, ist $\Pi\lambda$ der Wert des Integrals $\int y^\lambda e^{-y}\,\mathrm{d}y$ von $y=0$ bis $y=\infty$, sofern $\lambda+1$ eine positive Größe ist.

Setzt man allgemeiner $y=z^\alpha$, $\alpha\lambda+\alpha-1=\beta$, so geht $\int y^\lambda e^{-y}\,\mathrm{d}y$ in $\int\alpha z^\beta e^{-z^\alpha}\,\mathrm{d}z$ über; dieses wird also zwischen den Grenzen $z=0$ und $z=\infty$ durch $\Pi(\frac{\beta+1}\alpha-1)$ ausgedrückt

\noindent oder

Der Wert des Integrals $\int z^\beta e^{-z^\alpha}\,\mathrm{d}z$ von $z=0$ bis $z=\infty$ ist $=\frac{\Pi(\frac{\beta+1}\alpha-1)}\alpha
=\frac{\Pi\frac{\beta+1}\alpha}{\beta+1}$, sofern $\alpha$ und $\beta+1$ positive Größen sind (wenn beide negativ sind, wird das\SourceLeaf{161} Integral durch $-\frac{\Pi\frac{\beta+1}\alpha}{\beta+1}$ ausgedrückt). So ergibt sich beispielsweise für $\beta=0$, $\alpha=2$ der Wert des Integrals $\int e^{-zz}\,\mathrm{d}z$ zu $=\Pi\frac12=\frac12\BareRoot\pi$.

\SixthArticle{29}
Der berühmte \textsc{Euler} hat für die Summe der Logarithmen $\log1+\log2+\log3+\mathrm{usw.}+\log z$ die Reihe $(z+\frac12)\log z-z+\frac12\log2\pi
+\frac{\mathfrak A}{1.2z}-\frac{\mathfrak B}{3.4z^3}
+\frac{\mathfrak C}{5.6z^5}-\mathrm{usw.}$ hergeleitet, wobei $\mathfrak A=\frac16$, $\mathfrak B=\frac1{30}$, $\mathfrak C=\frac1{42}$ usw. die \textsc{Bernoulli}schen Zahlen sind. Durch diese Reihe wird also $\log\Pi z$ ausgedrückt; denn wenn dieser Schluss auf den ersten Blick auch auf ganzzahlige Werte beschränkt scheint, so zeigt eine genauere Betrachtung doch, dass die von \textsc{Euler} verwendete Entwicklung (Instit. Calc. Diff. Cap.\,\textsc{vi}. 159) zumindest auf positive gebrochene Werte mit demselben Recht wie auf ganze Zahlen angewandt werden kann: Sie setzt nämlich nur voraus, dass die in eine Reihe zu entwickelnde Funktion von $z$ so beschaffen ist, dass ihre Abnahme, wenn $z$ in $z-1$ übergeht, durch den Satz von \textsc{Taylor} dargestellt werden kann und dass zugleich diese Abnahme $=\log z$ ist. Die erste Bedingung beruht auf der \emph{Stetigkeit} der Funktion und gilt daher nicht für negative Werte von $z$, auf die jene Reihe folglich nicht ausgedehnt werden darf; die zweite Bedingung kommt dagegen der Funktion $\log\Pi z$ allgemein zu, ohne Beschränkung auf ganzzahlige Werte von $z$. Setzen wir also
\[
\begin{aligned}
{}\text{[58]}\quad &\log\Pi z=(z+\frac12)\log z-z+\frac12\log2\pi\\[4pt]
&+\frac{\mathfrak A}{1.2z}-\frac{\mathfrak B}{3.4z^3}
+\frac{\mathfrak C}{5.6z^5}-\frac{\mathfrak D}{7.8z^7}
+\mathrm{usw.}
\end{aligned}
\]
Da sich hieraus auch
\[
\begin{aligned}
&\log\Pi2z=(2z+\frac12)\log2z-2z+\frac12\log2\pi\\[4pt]
&+\frac{\mathfrak A}{1.2.2z}-\frac{\mathfrak B}{3.4.8z^3}
+\frac{\mathfrak C}{5.6.32z^5}-\frac{\mathfrak D}{7.8.128z^7}
+\mathrm{usw.}
\end{aligned}
\]
ergibt und nach Formel 57, wenn $n=2$ gesetzt wird,
\[
\log\Pi(z-\frac12)=\log\Pi2z-\log\Pi z
-(2z+\frac12)\log2+\frac12\log2\pi,\quad\text{ergibt sich}
\]
\[
\begin{aligned}
{}\text{[59]}\quad &\log\Pi(z-\frac12)=z\log z-z+\frac12\log2\pi\\[4pt]
&-\frac{\mathfrak A}{1.2.2z}+\frac{7\mathfrak B}{3.4.8z^3}
-\frac{31\mathfrak C}{5.6.32z^5}+\frac{127\mathfrak D}{7.8.128z^7}
-\mathrm{usw.}
\end{aligned}
\]
Diese beiden Reihen konvergieren für große Werte von $z$ anfangs recht rasch, sodass sich eine angenäherte Summe bequem und hinreichend genau ermitteln lässt. Es ist jedoch wohl zu beachten, dass sich für jeden gegebenen Wert von $z$, wie groß er auch sein mag, nur eine begrenzte Genauigkeit erreichen lässt, da die \textsc{Bernoulli}schen Zahlen eine hypergeometrische Reihe bilden und jene Reihen daher, wenn sie nur weit genug fortgesetzt werden, gewiss aus konvergenten in divergente übergehen. Im Übrigen lässt sich nicht leugnen, dass die Theorie solcher divergenten Reihen noch mit gewissen Schwierigkeiten belastet ist, über die wir vielleicht bei anderer Gelegenheit ausführlicher sprechen werden.
